\documentclass{article}
% ---------------------------------------------------------------
% Póster de revisión de literatura (A0 vertical)
% Finanzas I (055D), UNCP, 2026-II. Compilar con pdfLaTeX.
% ---------------------------------------------------------------
\usepackage[paperwidth=841mm,paperheight=1189mm,margin=0mm]{geometry}
\usepackage[utf8]{inputenc}
\usepackage[T1]{fontenc}
\IfFileExists{spanish.ldf}{\usepackage[spanish,es-noshorthands,es-nodecimaldot]{babel}}{}
\usepackage{helvet}
\renewcommand{\familydefault}{\sfdefault}
\usepackage[letterspace=120]{microtype}
\usepackage{anyfontsize}
\usepackage{graphicx}
\usepackage{xcolor}
\usepackage[most]{tcolorbox}
\usepackage{enumitem}
\usepackage{multicol}
\usepackage{xurl}
\usepackage[hidelinks]{hyperref}
\usepackage{qrcode}
\urlstyle{same}
\pagestyle{empty}
\setlength{\parindent}{0pt}
\hyphenpenalty=10000 \exhyphenpenalty=10000

% ---------------- Paleta ----------------
\definecolor{navy}{HTML}{8E0C11}
\definecolor{navyLight}{HTML}{B8121A}
\definecolor{gold}{HTML}{00A3E0}
\definecolor{acento}{HTML}{A6E1F7}
\definecolor{goldSoft}{HTML}{E8F6FC}
\definecolor{teal}{HTML}{353435}
\definecolor{tealSoft}{HTML}{F1EEEE}
\definecolor{ink}{HTML}{231F20}
\definecolor{muted}{HTML}{5E5757}
\definecolor{line}{HTML}{E5D7D7}
\definecolor{page}{HTML}{F6F1F1}
\definecolor{mist}{HTML}{F7D3D5}
\pagecolor{page}

% ---------------- Tamaños (ajustados por póster) ----------------
\newcommand{\bodysize}{\fontsize{26}{34.3}\selectfont}
\newcommand{\cardtitlesize}{\fontsize{30.7}{36.8}\selectfont}
\newcommand{\smallsize}{\fontsize{21.3}{27.7}\selectfont}

% ---------------- Estilos de tarjeta ----------------
\tcbset{
  card/.style={enhanced, sharp corners, colback=white, colframe=line, boxrule=0.9mm,
    left=14mm, right=14mm, top=12mm, bottom=12mm, fontupper=\bodysize\color{ink}\raggedright,
    overlay={\fill[#1] (frame.north west) rectangle ([yshift=-5mm]frame.north east);}},
  card/.default=navy,
}
\newcommand{\cardtitle}[2][navy]{{\color{#1}\cardtitlesize\bfseries\textls[110]{\MakeUppercase{#2}}}\par\vspace{7mm}}
\newcommand{\bul}{\raisebox{0.18em}{\color{gold}\rule{0.5em}{0.5em}}}
\setlist[itemize]{label=\bul, leftmargin=1.25em, labelsep=0.6em, itemsep=0.55em, topsep=0pt, parsep=0pt}
\newcommand{\pill}[3]{\tcbox[on line, enhanced, arc=7mm, boxrule=0.8mm, colback=#1, colframe=#2,
  left=7mm, right=7mm, top=3mm, bottom=3mm, fontupper=\fontsize{26}{30}\selectfont\bfseries]{#3}}

\begin{document}

% ======================= ENCABEZADO =======================
\begin{tcolorbox}[enhanced, sharp corners, colback=navy, boxrule=0pt, width=\paperwidth,
  left=34mm, right=34mm, top=20mm, bottom=20mm,
  overlay={\fill[gold] (frame.south west) rectangle ([yshift=5mm]frame.south east);
           \fill[navyLight] ([xshift=-160mm]frame.north east) -- (frame.north east) -- ([yshift=-160mm]frame.north east) -- cycle;}]
\begin{minipage}[c]{92mm}\centering\includegraphics[width=\linewidth]{Escudo_UNCP.png}\end{minipage}\hfill
\begin{minipage}[c]{555mm}\centering
{\color{acento}\fontsize{25}{30}\selectfont\bfseries\textls[160]{REVISIÓN DE LITERATURA \quad|\quad FINANZAS I \quad|\quad TEMA 61}\par}
\vspace{10mm}
{\color{white}\fontsize{78}{88.1}\selectfont\bfseries Factor investing in the Turkish equity market\par}
\vspace{8mm}
{\color{mist}\fontsize{38}{46.4}\selectfont\itshape Inversión por factores en el mercado bursátil turco\par}
\vspace{10mm}
{\color{white}\fontsize{29}{35}\selectfont\bfseries Umut Gökçen\par}
\vspace{9mm}
\pill{white}{white}{\color{navy}Q1}\hspace{6mm}\pill{navy}{white}{\color{white}\mdseries \textit{Borsa Istanbul Review}, 100884 (2026)}
\end{minipage}\hfill
\begin{minipage}[c]{92mm}\centering\includegraphics[width=\linewidth]{Logo_FE.png}\end{minipage}
\end{tcolorbox}
\vspace{-1mm}
\begin{tcolorbox}[enhanced, sharp corners, colback=goldSoft, boxrule=0pt, width=\paperwidth,
  left=34mm, right=34mm, top=8mm, bottom=8mm, halign=center]
{\color{navy}\fontsize{27}{36}\selectfont \textbf{Póster elaborado por:} Meglinho Sneyder Rojas Poma \enspace\textbar\enspace Código 2024200522B\par \fontsize{24}{32}\selectfont Escuela Profesional de Economía, Universidad Nacional del Centro del Perú \enspace\textbar\enspace Finanzas I (055D), 2026-II\par}
\end{tcolorbox}

\vfill
% ======================= CUERPO =======================
\hspace*{34mm}\begin{minipage}{773mm}
\begin{tcolorbox}[enhanced, sharp corners, colback=navy, boxrule=0pt, left=16mm, right=16mm, top=12mm, bottom=11mm,
  fontupper=\bodysize\raggedright, overlay={\fill[gold] (frame.north west) rectangle ([yshift=-5mm]frame.north east);}]
\begin{minipage}[t]{174.8mm}\raggedright{\color{acento}\smallsize\bfseries\textls[110]{\MakeUppercase{Mercado}}}\par\vspace{2mm}{\color{white}Borsa Istanbul (Turquía)}\end{minipage}\hfill\begin{minipage}[t]{174.8mm}\raggedright{\color{acento}\smallsize\bfseries\textls[110]{\MakeUppercase{Muestra}}}\par\vspace{2mm}{\color{white}171 a 513 empresas por año}\end{minipage}\hfill\begin{minipage}[t]{174.8mm}\raggedright{\color{acento}\smallsize\bfseries\textls[110]{\MakeUppercase{Periodo}}}\par\vspace{2mm}{\color{white}1994 a 2024}\end{minipage}\hfill\begin{minipage}[t]{174.8mm}\raggedright{\color{acento}\smallsize\bfseries\textls[110]{\MakeUppercase{Método}}}\par\vspace{2mm}{\color{white}Carteras por factores y alfas}\end{minipage}
\end{tcolorbox}
\end{minipage}

\vfill
\hspace*{34mm}\begin{minipage}{773mm}
\begin{tcbraster}[raster columns=3, raster equal height, raster column skip=16mm]
\begin{tcolorbox}[card]\cardtitle{Contexto}Muchos estudios prueban factores de valoración en mercados emergentes juntando varios países, lo que oculta cómo funcionan en cada uno. Para Turquía, la evidencia previa es contradictoria e incompleta.\end{tcolorbox}
\begin{tcolorbox}[card]\cardtitle{Problema}No se sabía con datos completos y actuales qué factores explican realmente los retornos de las acciones de Borsa Istanbul.\end{tcolorbox}
\begin{tcolorbox}[card]\cardtitle{Objetivo}Probar catorce predictores de retornos en Borsa Istanbul y evaluar si sirven para diseñar estrategias sistemáticas de inversión.\end{tcolorbox}
\end{tcbraster}
\end{minipage}

\vfill
\hspace*{34mm}\begin{minipage}{773mm}
\begin{tcolorbox}[card]\cardtitle{Metodología}
\begin{minipage}[t]{0.48\linewidth}
\begin{itemize}
\item Se usa una base sin sesgo de supervivencia.
\item Se forman carteras según cada predictor y se miden sus primas.
\item Se prueba la robustez frente a la inflación y el riesgo cambiario, y se arma una cartera multifactor.
\end{itemize}
\end{minipage}\hfill
\begin{minipage}[t]{0.48\linewidth}
\begin{itemize}
\item Se prueban 14 predictores de corte transversal con la base más completa y actual de Borsa Istanbul.
\item Las carteras se ponderan por valor y se reportan el quintil superior, el inferior y su diferencia.
\end{itemize}
\end{minipage}
\end{tcolorbox}
\end{minipage}

\vfill
\hspace*{34mm}\begin{minipage}{773mm}
\begin{tcolorbox}[card]\cardtitle{Tabla clave del estudio}
\begin{minipage}[c]{233.3mm}\centering
\fcolorbox{line}{white}{\includegraphics[width=\dimexpr233.3mm-2\fboxsep-2\fboxrule\relax]{figura.png}}
\end{minipage}\hfill
\begin{minipage}[c]{\dimexpr\linewidth-233.3mm-18mm\relax}\raggedright
{\color{navy}\bfseries Tabla 4.}\enspace{\itshape Carteras univariadas de factores comunes, 1995 a 2024: retorno en exceso y alfa anualizados (\%) del quintil inferior (Q1), superior (Q5) y de la diferencia Q5-Q1; estadístico $t$ entre paréntesis.}\par\vspace{8mm}
{\color{teal}\smallsize\bfseries\textls[110]{CÓMO LEERLA}}\par\vspace{4mm}
\begin{itemize}
\item Valor es el factor más fuerte: la diferencia Q5-Q1 rinde 18.57 \% ($t=3.39$) y su alfa es 17.31 \% ($t=3.50$).
\item La beta no paga prima: la diferencia Q5-Q1 es 2.22 \% con $t=0.44$.
\item Entre 2005 y 2024 el tamaño ($-12.91$ \%, $t=-2.31$) y valor (14.93 \%) son significativos.
\end{itemize}
\vspace{8mm}
{\color{muted}\smallsize Tabla tomada de Gökçen (2026).}
\end{minipage}
\end{tcolorbox}
\end{minipage}

\vfill
\hspace*{34mm}\begin{minipage}{773mm}
\begin{tcolorbox}[card=teal]\cardtitle[teal]{Resultados}
\begin{tcolorbox}[enhanced, sharp corners, colback=tealSoft, boxrule=0pt, left=14mm, right=14mm, top=10mm, bottom=11mm,
  borderline west={4mm}{0mm}{teal}, fontupper=\bodysize\color{ink}\raggedright]
{\color{teal}\smallsize\bfseries\textls[110]{HALLAZGO CENTRAL}}\par\vspace{4mm}
{\bfseries El valor es el factor más fuerte y robusto, con una prima superior al 17 \% anual.}
\end{tcolorbox}
\vspace{11mm}
\begin{minipage}[t]{0.48\linewidth}
\begin{itemize}
\item Tamaño, rentabilidad y momentum son significativos, pero solo desde 2005.
\item Las acciones de bajo riesgo rinden más que las de alto riesgo con cualquier medida de riesgo.
\item La inflación y el riesgo cambiario no anulan las primas.
\end{itemize}
\end{minipage}\hfill
\begin{minipage}[t]{0.48\linewidth}
\begin{itemize}
\item Una cartera multifactor solo de compra logra un alfa de 13 \% en dólares.
\item La relación entre beta y retorno es, en el mejor caso, plana; entre 2005 y 2024 el alfa de la beta es negativo.
\item La cartera multifactor tiene un ratio de Sharpe de 0.64 frente al índice de mercados emergentes.
\end{itemize}
\end{minipage}
\end{tcolorbox}
\end{minipage}

\vfill
\hspace*{34mm}\begin{minipage}{773mm}
\begin{tcbraster}[raster columns=3, raster equal height, raster column skip=16mm]
\begin{tcolorbox}[card, raster multicolumn=2]\cardtitle{Conclusiones}Estudiar un solo país da una visión más clara que juntar mercados. En Turquía, los factores de valor, tamaño, rentabilidad y momentum sirven para inversión sistemática.\end{tcolorbox}
\begin{tcolorbox}[card=gold, colback=goldSoft]\cardtitle{Aporte al tema 61}Advierte que en un mercado emergente más riesgo no siempre significa más retorno, lo que ayuda a leer con cuidado el costo de capital calculado con el CAPM para cada sector.\end{tcolorbox}
\end{tcbraster}
\end{minipage}

\vfill
% ======================= PIE =======================
\begin{tcolorbox}[enhanced, sharp corners, colback=navy, boxrule=0pt, width=\paperwidth,
  left=34mm, right=34mm, top=14mm, bottom=11mm,
  overlay={\fill[gold] (frame.north west) rectangle ([yshift=-5mm]frame.north east);}]
\begin{minipage}[c]{245mm}\raggedright
{\color{acento}\smallsize\bfseries\textls[110]{FUENTE DE DATOS DEL ESTUDIO}}\par\vspace{5mm}
{\color{white}\smallsize Base de datos de mercado y estados financieros de RASYONET para las acciones de Borsa Istanbul, de 1994 a 2024.}
\end{minipage}\hfill
\begin{minipage}[c]{385mm}\raggedright
{\color{acento}\smallsize\bfseries\textls[110]{REFERENCIA}}\par\vspace{5mm}
{\color{white}\smallsize\hangindent=1.2em\hangafter=1 Gökçen, U. (2026). Factor investing in the Turkish equity market. \textit{Borsa Istanbul Review}, Artículo 100884. \url{https://doi.org/10.1016/j.bir.2026.100884}\par}
\end{minipage}\hfill
\begin{minipage}[c]{100mm}\centering
\tcbox[enhanced, sharp corners, colback=white, boxrule=0pt, left=3mm, right=3mm, top=3mm, bottom=3mm]{\qrcode[hyperlink, height=46mm]{https://doi.org/10.1016/j.bir.2026.100884}}\par\vspace{4mm}
{\color{acento}\fontsize{17}{21}\selectfont\bfseries\textls[100]{ESCANEA Y LEE EL ARTÍCULO}}
\end{minipage}
\end{tcolorbox}
\end{document}
